\documentclass[11pt]{article}
% Diplomatic transcription macros (12 Nov 2026). Typed text is plain; the handwritten layer is blue,
% deletions red, typewriter strike-outs grey, doubtful readings in brackets with a question mark.
\usepackage[utf8]{inputenc}
\usepackage[T1]{fontenc}
\usepackage{lmodern}
\usepackage[a4paper,margin=2.2cm]{geometry}
\usepackage{amsmath,amssymb}
\usepackage{xcolor}
\usepackage{tikz-cd}
\usepackage{environ}
\setlength{\parindent}{0pt}
\setlength{\parskip}{4pt}
\title{R. Brown's translations of two letters to L. Breen, with his note and comments (cote 134-2)\\[4pt]\large Diplomatic transcription from the Montpellier scans}
\author{AI-transcribed by Claude Opus 5.5}
\date{12 November 2026}
\newcommand{\scanpage}[1]{\par\bigskip\noindent{\small\sffamily[scan #1]}\par}
\newcommand{\mspage}[1]{{\small\sffamily[p.\,#1]}}
\newcommand{\hand}[1]{\textcolor{blue}{#1}}
\let\over\relax
\newcommand{\over}[2]{\textcolor{gray}{#1}\textcolor{blue}{\ensuremath{^{\!\nearrow}}#2}}
\newcommand{\inserted}[2][]{\textcolor{blue}{$\langle$#2$\rangle$}}
\newcommand{\margin}[1]{\textcolor{blue}{\{\textsf{marg.:} #1\}}}
\newcommand{\deleted}[1]{\textcolor{red}{[\textsf{del.:} #1]}}
\newcommand{\xout}[1]{\textcolor{gray}{[\textsf{x:} #1]}}
\newcommand{\uline}[1]{\underline{#1}}
\newcommand{\uuline}[1]{\underline{\underline{#1}}}
\newcommand{\unsure}[1]{[#1?]}
\newcommand{\illegible}[1]{[\textsf{illeg.\ #1}]}
\newcommand{\diagram}[1]{\par{\centering #1\par}}
\newenvironment{diagramblock}{\par\begin{center}}{\end{center}\par}

\begin{document}
\maketitle

% R. Brown's translations of two letters to L. Breen, with his note and comments (cote 134-2): archive cote 134-2, site page 115 (PDF page 116).
% cover note. Draft compared: none.
% AI-transcribed by Claude Opus 5.5 (claude-opus-5-5, adaptive thinking, effort high), Message Batches API, batch brown_full (msgbatch_01Md3i5yK1PxkpJn5GHfEExE), 2026-11-13; not yet checked by eye.
\scanpage{134-2/115}
\uline{Note}: The following is an informal translation, prepared by R. Brown
(Bangor) (with assistance from L. Breen (Rennes) and also from J-L. Loday
(Strasbourg)) of the mathematical part of a handwritten letter from
A. Grothendieck to L. Breen, dated 17-19 July, 1975. I would like to
thank Professor Grothendieck for making available copies of this, and
the two related previous letters.

\mspage{\hand{116}}

\clearpage
% R. Brown's translations of two letters to L. Breen, with his note and comments (cote 134-2): archive cote 134-2, site page 117 (PDF page 118).
% comments. Draft compared: none.
% AI-transcribed by Claude Opus 5.5 (claude-opus-5-5, adaptive thinking, effort high), Message Batches API, batch brown_full (msgbatch_01Md3i5yK1PxkpJn5GHfEExE), 2026-11-13; not yet checked by eye.
\scanpage{134-2/117}
\uline{COMMENTS}

1. The topic of duality has progressed somewhat since 1975, see for
example

L. Begueri \hand{Memoir SMF}
\margin{demander réf. : Breen.}

M. Bester \hand{Thesis, U of Michigan}

(These references are due to L. Breen.)

2. The work of Mme Sinh referred to on p. 2 is as follows:

H.X. Sinh, \textit{Gr-catégories}, Thèse Doctorat d'\over{E}{É}tat Université Paris VII
(1975).

See also

N.\deleted{\unsure{S}} \inserted{Saavedra} Rivano, \textit{Catégories tanakiennes}, Springer Lecture
Notes in Mathematics 265 (1972).

Related to these notions is that of +-groupoid (=coherent
symmetric monoidal groupoid with strict associativity and strict
zero object) whose theory is developed also in:

M. Takeuchi, 'On Villamayor and Zelinsky's long exact sequence',
Mem. Amer. Math. Soc. 33 (1981) No 249.

3. The description of groups in (1-cat) given on p. 2-3 is in effect an
equivalence between G-groupoids (= group objects in the category of
groupoids) and crossed modules. This is described also in:

R. Brown and C.B. Spencer, 'G-groupoids, crossed modules and the
fundamental groupoid of a topological group', Proc. Kon. Akad.
v. Weten. 79(1976), 296 - 302.
and is believed to have been known to Verdier in 1965.

4. The equivalence mentioned in the last comment has been generalised to

\mspage{118}

\clearpage
% R. Brown's translations of two letters to L. Breen, with his note and comments (cote 134-2): archive cote 134-2, site page 118 (PDF page 119).
% comments. Draft compared: none.
% AI-transcribed by Claude Opus 5.5 (claude-opus-5-5, adaptive thinking, effort high), Message Batches API, batch pilotnext (msgbatch_01A9sdu8SYq23bEtqX9tUncQ), 2026-11-13; not yet checked by eye.
\scanpage{134-2/118}
\mspage{-2-}

equivalences between six categories as set out in the diagram,

\begin{diagramblock}\begin{tikzcd}[column sep=large, row sep=large]
\text{poly -T-complexes} \arrow[r, leftrightarrow, "(f)"] \arrow[d, leftrightarrow, "(f)"'] & \text{cubical T-complexes} \arrow[r, leftrightarrow, "(c)"] & \omega\text{-groupoids} \arrow[d, leftrightarrow, "(d)"] \arrow[dl, leftrightarrow, "(a)"] \\
\text{simplicial T-complexes} \arrow[r, leftrightarrow, "(e)"'] & \text{crossed complexes} \arrow[r, leftrightarrow, "(d)"'] & \infty\text{-groupoids}
\end{tikzcd}\end{diagramblock}

The categories are defined and the equivalences proved in:

a) R. Brown and P. J. Higgins, 'The algebra of cubes',
J. Pure Appl. Alg. 21(1981) 233 - 260.

b) R. Brown and P. J. Higgins, 'Colimit theorems for relative
homotopy groups', \textit{ibid.} 22(1981) 11 - 41.

c) R. Brown and P. J. Higgins, 'The equivalence of $\omega$-groupoids and
cubical T-complexes', Cah. Top. Géom. Diff. (3$^{\text{e}}$ Coll. sur les
catégories, dédié a Charles Ehresmann) 22 (1981), 349 - 369.

d) R. Brown and P. J. Higgins, 'The equivalence of $\infty$-groupoids and
crossed complexes' \textit{ibid,} 370 - 386.

e) N. K. Ashley, \textit{Simplicial T-complexes,} Ph.D. thesis, University
of Wales (1978),(to be published in an issue of Esquisses
Mathématiques, Amiens).

f) D. W. Jones, Ph.D. thesis in preparation.

Applications of the equivalence (a) are given in reference (b)
above.

5. However the notion of $\infty$-groupoid referred to in (4) is not adequate
for the purposes considered by Grothendieck since, roughly speaking,
crossed complexes describe the homotopy types only of spaces X
which fibre over a K($\pi$,1) with fibre a product of Eilenberg-
MacLane spaces.

6. A notion of multiple groupoid called an \textit{n-cat-group} has been described

\mspage{\hand{119}}

\clearpage
% R. Brown's translations of two letters to L. Breen, with his note and comments (cote 134-2): archive cote 134-2, site page 119 (PDF page 120).
% comments. Draft compared: none.
% AI-transcribed by Claude Opus 5.5 (claude-opus-5-5, adaptive thinking, effort high), Message Batches API, batch brown_full (msgbatch_01Md3i5yK1PxkpJn5GHfEExE), 2026-11-13; not yet checked by eye.
\scanpage{134-2/119}
\mspage{-3-}

in:

J.-L. Loday, 'Spaces with finitely many homotopy groups', J. Pure
Appl. Alg.  24(1982) 179 - 202.
The n-cat-groups (which may be described as n-fold categories in
the category of groups) are adequate for the description of
truncated homotopy types.  This is Loday's main result.

7. An analysis of the structure of non-abelian chain complex on a
length two normalised chain complex of a simplicial group is
given in:
D. Conduché, 'Modules croisés généralisés de longeur 2' (in
preparation).

8. By following through some of the equivalences in (4) above, we
can associate to an $\infty$-groupoid G \quad a simplicial T-complex NG, whose
underlying simplicial set could be called the \textit{nerve} of G.  A
direct description of NG is that $(\mathrm{NG})_{n}$ is the set of
$\infty$-groupoid maps  to G from $\sigma\underset{\approx}{\Delta}^{n}$, the homotopy $\infty$-groupoid of the
standard n-simplex $\Delta^{n}$ (i.e. $\sigma\underset{\approx}{\Delta}^{n}$ has associated crossed complex $\pi\underset{\approx}{\Delta}^{n}$,
the homotopy crossed complex of $\underset{\sim}{\Delta}^{n}$).
Now a definition of $\infty$-category is also given in (4)(d) above,
but the corresponding functor nerve: $\infty$-categories $\longrightarrow$ simplicial
sets is still  unknown.

9. An indication of the use of the ideas of (4) in non-abelian
situations is given in:

a) R. Brown and P.J. Higgins, 'Crossed complexes and non-%
abelian extensions', Int. Conf. on Category Theory, Gummersbach
(1981),  Springer L.N.M. (to appear).

b) R. Brown, 'Non-abelian cohomology and the homotopy classification
of maps', Proc. Conf, \textit{Méthodes d'algèbre homotopique en topologie},
Marseilles - Luminy (1982).  Astérixes (to appear)

\mspage{\hand{120}}

\clearpage
% R. Brown's translations of two letters to L. Breen, with his note and comments (cote 134-2): archive cote 134-2, site page 120 (PDF page 121).
% comments. Draft compared: none.
% AI-transcribed by Claude Opus 5.5 (claude-opus-5-5, adaptive thinking, effort high), Message Batches API, batch brown_full (msgbatch_01Md3i5yK1PxkpJn5GHfEExE), 2026-11-13; not yet checked by eye.
\scanpage{134-2/120}
\mspage{-4-}

The latter paper uses the equivalence between crossed complexes
and simplicial T-complexes (which generalises the Dold-Kan
theorem on the equivalence of chain complexes and simplicial
abelian groups) to give a notion of nerve, and so classifying
space BC of a crossed complex C. The classical Eilenberg-%
MacLane bijection
\[ [X, K(A, n)] \cong H^{n} (X, A) \]
for a CW-complex X, is then generalised to a bijection
\[ [X, BC] \cong H^{0} (X, C) \]
However, as pointed out in (5), these results are still ``too
abelian'', and much more work is needed to realise Grothendieck's
programme.

10. The theory of tubular neighbourhoods has been developed in papers
by D.A. Cox.

'Algebraic tubular neighbourhoods I, II', \unsure{Math,|Math.} Scand. 42 (1978),
211 - 228, 229 - 242.

Another of his papers seems relevant:

'Homotopy theory of simplicial schemes', Compositio Math. 39(1979)
263 - 296.

R. BROWN

School of Mathematics and Computer Science,
University College of North Wales,
Bangor,
Gwynedd. \hand{U.K.}

November, 1982

\margin{11 A van Kampen Theorem for 2-cat-groups has now
been proved by Brown and Loday.}

\mspage{\hand{121}}

\clearpage
% R. Brown's translations of two letters to L. Breen, with his note and comments (cote 134-2): archive cote 134-2, site page 122 (PDF page 123).
% translation of the letter of 17/2/1975. Draft compared: none.
% AI-transcribed by Claude Opus 5.5 (claude-opus-5-5, adaptive thinking, effort high), Message Batches API, batch brown_full (msgbatch_01Md3i5yK1PxkpJn5GHfEExE), 2026-11-13; not yet checked by eye.
\scanpage{134-2/122}
\uline{Translation of a letter of A. Grothendieck to L. Breen 17/2/1975}

Dear Larry,

Here is an afterthought to ``une lettre-fleuve'' on the yoga of homotopy.
As you doubtless know, to a topos $X$ one associates canonically a pro-simplicial
set, and so in a convenient sense a ``pro-homotopy type''. When $X$ is ``locally
homotopically trivial'', the associated pro-object is essentially constant as
a pro-object in the homotopy category, and so $X$ defines, in the usual homotopy
category, an object which is the ``homotopy type''. Similarly, if $X$ is
``locally homotopically trivial in $\dim \leq n$'' , it defines an ordinary homotopy
type, but ``truncated in $\dim \leq n$'' - this is a familiar construction for $n = 0$
or $1$ , even among those like me who know hardly any homotopy theory!

These constructions are functorial in $X$ . Moreover, if $f : X \to Y$ is
a morphism of toposes, Artin-Mazur have given a \textit{cohomological} condition which \margin{topoi}
is necessary and sufficient for $f$ to be a ``homotopy equivalence in $\dim \leq n$'':
it is that $H^{i}(Y, F) \xrightarrow{\cong} H^{i}(X, f^{*}(F))$ for $i \leq n$ , and all \textit{locally constant}
sheaves of groups $F$ on $Y$ , allowing for $i \leq 1$ that $F$ be non-commutative.
This criterion, in terms of ``locally constant'' n-gerbes $F$ on $Y$ , can be
interpreted as the condition that $F(Y) \to F(X)$ is an n-equivalence for all
such $F$ and $i \leq n$ . It is certainly true that this is equivalent to the
following criterion:

(A) For every ``locally constant'' n-champ $F$ on $Y$ , the n-functor \margin{\uline{Stack}}
$F(Y) \to f^{*}(F)(X)$ is an n-equivalence,

or again

(B) The n-functor $F \to f^{*}(F)$ which sends the n-category of locally
constant (n-1)-champs in $Y$ to that of locally constant (n-1)-champs on
$X$ , is an n-equivalence.

In other terms, the constructions on a topos $X$ which one can make in
terms of (n-1)-champs which are \textit{locally} constant, depend only on its ``n-%
truncated pro-homotopy type'', and define it. In the case where $X$ is locally
homotopically trivial in $\dim \leq n$ , and so defines a n-truncated ordinary
homotopy type, one can interpret these last as an n-groupoid $C_{n}$ , (defined
up to n-equivalence). In terms of these

(C) The (n-1)-champs on $X$ should be able to be identified with the
n-functors from the category $C_{n}$ in the n-category ((n-1)-Cat) of all
(n-1)-categories.

In the case $n = 1$ , this is nothing other than the Poincaré theory of
the classification of coverings of $X$ in terms of the``fundamental groupoid''
$C$ , of $X$ . By extension, $C_{n}$ merits the name \textit{fundamental n-groupoid} of $X$ ,
which I propose to write $\Pi_{n}(X)$ . Knowledge of this includes knowledge of
the $\pi_{i}(X)$ $(0 \leq i \leq n)$ and the Postnikoff invariants of all orders up to
$H^{n+1}(\Pi_{n-1}(X), \pi_{n})$ .

In the case of an arbitrary topos, not necessarily locally homotopically
trivial in $\dim \leq n$ , one hopes to be able to interpret the (n-1)-champs which
are locally constant on $X$ in terms of a $\Pi_{n}(X)$ which will be a pro-n-groupoid.
This has been done, more or less, for $n = 1$ (at least for connected $X$) ;
the case where $X$ is the étale topos of a scheme is treated extensively in
SGA3, in relation to the classification of tori on an arbitrary base.

\mspage{123}

\clearpage
% R. Brown's translations of two letters to L. Breen, with his note and comments (cote 134-2): archive cote 134-2, site page 123 (PDF page 124).
% translation of the letter of 17/2/1975. Draft compared: none.
% AI-transcribed by Claude Opus 5.5 (claude-opus-5-5, adaptive thinking, effort high), Message Batches API, batch brown_full (msgbatch_01Md3i5yK1PxkpJn5GHfEExE), 2026-11-13; not yet checked by eye.
\scanpage{134-2/123}
\mspage{- 2 -}

In the case $n = 1$ , one knows that one can recover (up to equivalence)
the 1-groupoid $C_1$ from the 1-category $\mathit{Hom}(C_1, (\mathrm{Ens}))$ of the functors into
(Ens) = (0 - Cat) (i.e. the ``local systems'' on $C_1$) - which is a topos, called
``multigalois'' - like the category of ``fibred functors'' on the above topos,
i.e. the opposite category to the category of points of this topos (which
is none other than the \textit{classifying topos} of $C_1$) . To make precise for
arbitrary $n$ the way in which the homotopy n-type of a topos $X$ (supposed
for simplicity to be locally homotopically trivial in dim $\leq n$ ) i.e. its
fundamental n-groupoid $C_n$ , can be expressed in terms of the n-category of
``local (n-1)-systems on $X$'' i.e. of the locally constant (n-1)-champs on $X$ ,
and to elucidate completely the hypothetical statement (B) above, it is
necessary to make explicit how an n-groupoid $C_n$ can be recovered, up to
n-equivalence, from the knowledge of the n-category
\[ \mathbb{C}_n = (n\text{-}\mathit{Hom})(C_n , ((n-1)\text{-}\mathrm{Cat})) \]
of local (n-1)-systems on $C_n$ . One would like to say that $C_n$ is the category
of ``fibred n-functors'' on $\mathbb{C}_n$ , i.e. of n-functors $\mathbb{C}_n \rightarrow ((n-1)\text{-}\mathrm{Cat})$ having
certain exactness properties (for $n = 1$ ; this is the condition of being
the inverse image functor for a morphism of toposes, i.e. to commute with
arbitrary $\underrightarrow{\ell im}$ and with finite $\underleftarrow{\ell im}$ ....). It is this which makes real
the fear, expressed in my preceding letter, that one ends by falling upon
the notion of n-topos and of morphisms of these! $\mathbb{C}_n$ will be an n-topos,
(called the ``classifying n-topos'' of the n-groupoid $C_n$) , ((n-1)-Cat) will
be the n-topos of points, and $C_n$ will be interpreted modulo n-equivalence
as the n-category of ``n-points'' of the classifying n-topos $\mathbb{C}_n$ . Brr!

If one hopes to be able to define a good old classifying (1)-topos for
an n-groupoid $C_n$ , as solution of a universal problem, I can see only how
\margin{\unsure{this} seems \unsure{suspect}} to recover the following universal problem: for every topos $T$ , consider
$\mathit{Hom}(\Pi_n(T), C_n)$ . This is an n-category, but take from it the truncated
1-category $\tau_1 \mathit{Hom}(\Pi_n(T), C_n)$ . For variable $T$ , one wants to 2-represent
the contravariant 2-functor $(\mathrm{Top})^{\circ} \rightarrow (1\text{-}\mathrm{Cat})$ by a classifying topos
$\mathbb{B} = \mathbb{B}_{C_n}$ , and then to find a 2-universal $\Pi_n(\mathbb{B}) \xrightarrow{\ \phi\ } C_n$ in the sense that
for all $T$ , the functor
\[ \mathit{Hom}_{(\mathrm{Top})}(T, \mathbb{B}) \longrightarrow \tau_1 \mathit{Hom}(\Pi_n(T), C_n) \]
\[ u \longmapsto \phi_0 \Pi_n(u) \]
is an equivalence. For $n = 1$ one knows that the usual classifying topos
of $C_1$ does the job, but for $n = 2$ already, I doubt that this universal
problem has a solution. This is perhaps related to the fact that the
``van Kampen Theorem'', which one can express by saying that the 2-functor
$T \longmapsto \Pi_1(T)$ of locally 1-connected toposes to groupoids transforms (up to
1-equivalence) amalgamated sums to amalgamated sums (and more generally commutes
with inductive 2-limits), is doubtless no longer true for $\Pi_2(T)$ . Thus,
if $T$ is a topological space which is the union of two closed sets, $T_1$ and
$T_2$ , it is doubtless not true that giving a locally constant 1-champ
on $T$ ``is equivalent to'' giving a locally constant 1-champ $F_i$ on $T_i$

\mspage{124}

\clearpage
% R. Brown's translations of two letters to L. Breen, with his note and comments (cote 134-2): archive cote 134-2, site page 124 (PDF page 125).
% translation of the letter of 17/2/1975. Draft compared: none.
% AI-transcribed by Claude Opus 5.5 (claude-opus-5-5, adaptive thinking, effort high), Message Batches API, batch pilotnext (msgbatch_01A9sdu8SYq23bEtqX9tUncQ), 2026-11-13; not yet checked by eye.
\scanpage{134-2/124}
\mspage{- 3 -}

(i = 1, 2) and an equivalence between the restrictions of $F_1$ and $F_2$
to $T_1 \cap T_2$ (while the analogous statement in terms of 0-champs, i.e. for
coverings, is evidently correct).

The statement (B) above makes it clear how to give explicitly the
cohomology of an n-groupoid $C_n$ . If $C_n = \Pi_n(X)$ , and if $F$ is a locally
constant (n-1)-champ on $X$ , and $e^{X}_{n-1}$ is the ``final'' (n-1)-champ, one has
an (n-1)-equivalence of (n-1)-categories
\[ \Gamma_{\unsure{x|\times}}(F) = F(X) \simeq \mathit{Hom}(e^{X}_{n-1}, F) \]
which shows that the functor $\Gamma_X$ ``integration on $X$'' for locally constant
(n-1)-champs, which includes the (non-commutative) locally constant cohomology
of $X$ in dim $\leq$n-1, can be interpreted in terms of ``local (n-1)-systems''
on the fundamental groupoid as an
$\mathit{Hom}(e^{(C_n)}_{n-1} , F)$ where now $F$ is interpreted as an n-functor
\[ C_n \xrightarrow{\ F\ } ((n-1)\text{-Cat}) \]
and $e^{(C_n)}_{n-1}$ is the constant n-functor on $\unsure{\underline{C}_n|C_n}$ , with value the final (n-1)-
category.

To interpret this in cohomology notation, it is necessary for me to add,
as ``apology'' to the preceding letter, the explicit interpretation of the
non-commutative cohomology on a topos $X$ , in terms of integration of n-
champs on $X$ . If $F$ is a strict Picard n-champ on $X$ , then it is defined
by a complex $L^{\circ}$ of cochains on $X$
\[ 0 \to L^0 \to L^1 \to L^2 \to \ldots \to L^n \to 0 , \]
concentrated in degrees $0 \leq i \leq n$ (defined uniquely up to isomorphism in
the derived category of Ab(X)). That said, the $H^i(X, L^{\circ})$ (hypercohomology)
for $0 \leq i \leq n$ can be interpreted as $H^i(X, L^{\circ}) = \pi_{n-i}\Gamma_X(F)$ . If one is
interested in all the $H^i$ (not just for $i \leq n$) one must, for all $N \geq n$ ,
regard $L^{\circ}$ as a complex concentrated in degrees $0 \leq i \leq N$ (by prolongation
of $L^{\circ}$ by 0 to the right). The corresponding strict Picard n-champ is no
longer $F$ but $\mathbb{C}^{N-n}F$ , where $\mathbb{C}$ is the ``classifying space'' functor, interpreted
on strict Picard n-categories as the operation consisting of ``translating''
the i-objects to (i+1)-objects, and adjoining a unique 0-object; this extends
one hopes, in ``an obvious way'', to n-champs, so as to commute with the operation
of taking the inverse image of an n-champs. One has then for $i \leq N$
\[ H^i(X, L^{\circ}) = \pi_{N-i}\Gamma_X(\mathbb{C}^{N-n}F) \qquad i \leq N . \]
Given this, it is necessary to put, for all strict Picard n-champs $F$ on $X$ ,
\[ H^i(X, F) = \pi_{N-i}\Gamma_X(\mathbb{C}^{N-n}F) \quad \text{if} \quad N \leq i, n \]
which does not depend on the choice of integer $N \geq \sup(i, n)$ .

[\uline{N.B.} One has a canonical morphism of (n-1)-groupoids
\[ \mathbb{C}(\Gamma_X F) \longrightarrow \Gamma_X(\mathbb{C}F) , \]
as the obvious constructions in terms of cochains show, and one sees in the

\mspage{\hand{125}}

\clearpage
% R. Brown's translations of two letters to L. Breen, with his note and comments (cote 134-2): archive cote 134-2, site page 125 (PDF page 126).
% translation of the letter of 17/2/1975. Draft compared: none.
% AI-transcribed by Claude Opus 5.5 (claude-opus-5-5, adaptive thinking, effort high), Message Batches API, batch brown_full (msgbatch_01Md3i5yK1PxkpJn5GHfEExE), 2026-11-13; not yet checked by eye.
\scanpage{134-2/125}
\mspage{- 4 -}

same way that this induces isomorphisms on $\pi_i$ for $1 \leq i \leq n + 1$ .]

[\uline{N.B}. One sees by the way that for $F$ an n-champ of groupoids on $X$ ,
if one restricts to defining the $H^i(X, F)$ for $0 \leq i \leq n$ , one has no need
of a Picard structure on $F$ , as it is sufficient to put
\[ H^i(X, F) = \pi_{n-i}(\Gamma_X(F)) \qquad 0 \leq i \leq n \ . \]
If on the other hand $F$ is an n-Gr-champ (i.e. $F$ has the structure of a
composition law $F \times F \rightarrow F$ with the usual formal properties of a group) the
``classifying (n+1)-champ'' is defined, and one can define $H^i(X, F)$ for
$0 \leq i \leq n + 1$ by
\[ H^i(X, F) = \pi_{n+1-i}(\Gamma_X(\mathbb{C}F)) \]
in particular
\[ H^{n+1}(X, F) = \pi_0(\Gamma_X(\mathbb{C}F)) \]
\[ \hspace{6em} = \text{equivalence classes of sections of } \mathbb{C}F \ . \]
But one can form $C\mathbb{C}F = \mathbb{C}^2F$ and define $H^{n+2}(X, F)$ , it seems \textit{only} if $\mathbb{C}F$
is itself a Gr-(n+1)-champ, which is without doubt the case only if $F$ is
a strict Picard n-champ ....

Let us now come to the case where $F$ is a \textit{locally constant} n-champ on
$X$ , and so is defined by an (n+1)-functor
\[ C_{n+1} \xrightarrow{F} \text{(strict Picard n-Cat)} \ . \]

Then, putting for $0 \leq i \leq n$
\[ H^i(C_{n+1}, F) = \pi_{n-i}(\mathit{Hom}(e_n^{(C_{n+1})}, F)) \ , \]
``one knows it fails'', as one has a canonical isomorphism
\[ H^i(C_{n+1}, F) \cong H^i(X, F) \]

valid in effect without Picard structure on $F$ ... It is thus necessary
for all $i$ and for every $\infty$-groupoid $C$ and every (n+1)-functor
\[ C \xrightarrow{F} \text{(strict Picard n-Cat)} \ , \]
to define
\[ H^i(C, F) = \pi_{N-i}\mathit{Hom}(e_N^C, \mathbb{C}^{N-n}F) \]

where one chooses $N \geq \sup(i, n)$ , If $F$ has only a Gr-structure (not
necessarily Picard) one can define the $H^i(C, F)$ for $i \leq n + 1$ by
\[ H^i(C, F) = \pi_{n+1-i}\mathit{Hom}(e^C_{n+1}, \mathbb{C}F) \ . \]
In the case $C = C_{n+1} = \Pi_{n+1}(X)$ , it must still be true (by virtue of (A)
above), that this set is canonically isomorphic to $H^{n+1}(X, F) = \pi_0\Gamma_{\unsure{x|\times}}(\mathbb{C}F)$
(this is true and very easy for $n = 0$). Can one describe the arrow
between the two sides of
\[ H^{n+1}(X, F) \simeq H^{n+1}(\Pi_{n+1}(X), F) \ ? \]

If one wishes to make (A) and (B) explicit again, in terms of the yoga
(C), one comes to the following situation:

\mspage{\hand{126}}

\clearpage
% R. Brown's translations of two letters to L. Breen, with his note and comments (cote 134-2): archive cote 134-2, site page 126 (PDF page 127).
% translation of the letter of 17/2/1975. Draft compared: none.
% AI-transcribed by Claude Opus 5.5 (claude-opus-5-5, adaptive thinking, effort high), Message Batches API, batch brown_full (msgbatch_01Md3i5yK1PxkpJn5GHfEExE), 2026-11-13; not yet checked by eye.
\scanpage{134-2/126}
\mspage{- 5 -}

One has an (n+1)-functor between (n+1)-groupoids
\[ f_{n+1} : C_{n+1} \longrightarrow D_{n+1} \]
which induces by truncation an n-functor
\[ f_n : C_n \longrightarrow D_n \, . \]
One must then have:

(A') $f_n$ is an n-equivalence if and only if the n-functor
\[ f_n^{*} : \mathit{Hom}(D_n, ((n-1)\text{-Cat})) \longrightarrow \mathit{Hom}(C_n, ((n-1)\text{-Cat})) \]
which sends the local (n-1)-systems on $D_n$ (or, equally, on $D_{n+1}$)
to the local (n-1)systems on $C_n$ , is an n-equivalence.

(B') $f_n$ is an n-equivalence if and only if for every local n-system $F$
on $D_{n+1}$ ,
\[ F : D_{n+1} \longrightarrow (n\text{-Cat}) \ , \]
the n-functor induced by $f_{n+1}$
\begin{diagramblock}
\begin{tikzcd}
\mathit{Hom}(e_n^{D_{n+1}}, F) \arrow[r] \arrow[d, equal, "\text{def.}"] & \mathit{Hom}(e_n^{C_{n+1}}, f^{*}_{n+1} F) \arrow[d, equal, "\text{def.}"] \\
\Gamma_{D_{n+1}}(F) & \Gamma_{C_{n+1}}(F)
\end{tikzcd}
\end{diagramblock}
is an n-equivalence.

The construction of the cohomology of a topos in terms of integration
of champs makes no appeal at all to complexes of abelian sheaves and still
less to the technique of injective resolutions. One has the impression that
in this spirit, \textit{via} the definition (which remains to be made explicit!) of
n-champs, it is all related above all to the ``\v{C}echist'' calculations in terms
of hypercoverings. Now these last are written with the help of a small dose
of semi-simplicial algebra. I do not know if a theory of champs and of
operations on them can be written \textit{without} ever using semi-simplicial algebra.
If yes, there would be essentially three distinct approaches for constructing
the cohomology of a topos:

(a) viewpoint of complexes of sheaves, injective resolutions, derived
categories (\textit{commutative homological algebra})

(b) viewpoint \v{C}echist or semi-simplicial (\textit{homotopical algebra})

(c) viewpoint of n-champs (categorical algebra, or \textit{non-commutative homological
algebra}).

In (a) one ``resolves'' the coefficients, in (b) one resolves the base
space (or topos), and in (c) it appears one resolves neither the one nor the
other.

Very cordially,

Alexandre.

\mspage{\hand{127}}

\clearpage
% R. Brown's translations of two letters to L. Breen, with his note and comments (cote 134-2): archive cote 134-2, site page 128 (PDF page 129).
% translation of the letter of 17-19 July 1975. Draft compared: ps1.tex, lines 1380-1404.
% AI-transcribed by Claude Opus 5.5 (claude-opus-5-5, adaptive thinking, effort high), Message Batches API, batch brown_full (msgbatch_01Md3i5yK1PxkpJn5GHfEExE), 2026-11-13; not yet checked by eye.
\scanpage{134-2/128}
% NOTE: a white patch on the scan, after "from" in the first title line, hides part of the line
\uline{Translation of an extract from}

% NOTE: a white patch on the scan, before "a Letter" in the second title line, hides part of the line
\uline{a Letter from A. Grothendieck to L. Breen}

\uline{17-19 July, 1975}

I am happy to finish by receiving an echo to my long letter
and even a beginning to a constructive approach to a theorem of
the type I envisaged.  The construction which you propose for the
notion of a non-strict n-category, and of the functor nerve, has
certainly the merit of existing, and of being a first precise
approach, but otherwise can be  subject to some evident criticism:
it is very technical, hardly intuitive (yet at the level of
1-cat, etc, for 2-cat everything  is so clear ``you just follow
your nose ...'')  And finally the absence of a definition of a
functor sending (semi-)simplicial sets to n-groupoids.  This
functor corresponds to a geometric intuition so clear that a theory
which does not include it seems to me to be useless.  Perhaps
in trying to write down (like a sort of list of Christmas presents!)
in a complete and explicit enough way the notions which one would
like to have at ones disposal, and the relations (functor,
equivalence, etc) one will arrive finally at a kind of axiomatic
description sufficiently complete to be able to give the key to
an explicit \uline{ad} \uline{hoc} construction, or should permit maybe to enunciate
and prove a theorem of existence and uniqueness for a theory of
the required type.

Otherwise, not having understood the idea of Segal in your
last letter (which I have generously sent to Illusie ...), I do
not see how you define the Picard n-categories - but this matters
little.  As far as ``strict'' Picard n-categories are concerned,
what I ask of them is that they finally form an (n+1)-category

\mspage{\hand{129}}

\clearpage
% R. Brown's translations of two letters to L. Breen, with his note and comments (cote 134-2): archive cote 134-2, site page 129 (PDF page 130).
% translation of the letter of 17-19 July 1975. Draft compared: ps1.tex, lines 1385-1394.
% AI-transcribed by Claude Opus 5.5 (claude-opus-5-5, adaptive thinking, effort high), Message Batches API, batch brown_full (msgbatch_01Md3i5yK1PxkpJn5GHfEExE), 2026-11-13; not yet checked by eye.
\scanpage{134-2/129}
\mspage{- 2 -}

(n+1)-equivalent to that of complexes of chains of length $n$ .
Agreed?  I thank you for having rectified in my mind a big blunder,
because of my great ignorance of algebraic topology and homotopy -
I was in fact of the impression that $H$-spaces satisfying conditions
of associativity and commutativity strict enough (say equivalent
to an $\Omega^{i}X$ with $i$ arbitrarily large) correspond to commutative
topological groups (and so inspiring in me several analogies ...).
Thus I have intuitive agreement with your observations on p.5.

On the other hand, I am always intrigued by the following
question: is there an analogue of the ``tapis'' of Dold-Puppe for
semi-simplicial groups (not necessarily commutative) and what form
should it take?  To say it truly I am considering the yoga
\[ \text{simplicial sets } \hand{$\longleftrightarrow$} \ \infty\text{-groupoids} \]
as being essentially the ultimate ``set theoretic'' version of
Dold-Puppe, which I would deduce by making explicit solely the
fact that the abelian groups in  ($\infty$-cat) are ``nothing else'' than
the complexes of chains in  (Ab) .  It is necessary then to
determine what should be the groups in ($\infty$-cat). I can tell you what
these are in (1-cat) - this will be discussed at length in the book
of Mme. Sinh, I think, under the title of the chapter ``\uline{strict}
Gr-categories'' (i.e. the isomorphisms of associativity, unity and
inverse  $X.X^{-1} \simeq 1$  are \uline{identities}).  One can make explicit how
(\uline{via} the fact that a Gr-category is Gr-equivalent to a strict
Gr-category)      the calculation with the Gr-categories reduces
to a very algebraic calculation with the \uline{strict} Gr-categories,
by a kind of ``calculus of fractions'' (by choice, left or right)
of the type which ends in giving the construction of derived

\mspage{\hand{130}}

\clearpage
% R. Brown's translations of two letters to L. Breen, with his note and comments (cote 134-2): archive cote 134-2, site page 130 (PDF page 131).
% translation of the letter of 17-19 July 1975. Draft compared: ps1.tex, lines 1394-1418.
% AI-transcribed by Claude Opus 5.5 (claude-opus-5-5, adaptive thinking, effort high), Message Batches API, batch pilotnext (msgbatch_01A9sdu8SYq23bEtqX9tUncQ), 2026-11-13; not yet checked by eye.
\scanpage{134-2/130}
\mspage{- 3 -}

categories. Whatever it should be, here is the explicit formulation
of the structure (groups in (1-cat)) in terms of the theory of
groups (1-categories in (Gr)). The structure is described by a
quadruplet $(L_1, L_0, d,\theta)$ with
\[ L_1 \xrightarrow{\ d\ } L_0 \]
a homomorphism of ordinary groups,
\[ \theta : L_0 \longrightarrow \mathrm{Aut}_{Gr}(L_1) \]
an operation of $L_0$ on $L_1$ , with the following two axioms:
(a) $d$ commutes with the operation of $L_0$ , in the
sense that
\[ d(\theta(x_0).x_1) = \mathrm{int}(x_0)\ d(x_1) \]
(b) \[ \theta(d(x_1)) = \mathrm{int}(x_1) \ . \]
These properties imply that $\mathrm{Im}\, d$ is normal in $L_0$ , (hence
$\mathrm{Coker}\ d = \pi_0$ is defined) and $\pi_1 = \mathrm{Ker}\ d$ is central in $L_1$ ,
and finally that $L_0$ operates on $L_1$ leaving $\pi_1$ invariant,
and it operates \uline{via} $\pi_0$ . The principal cohomological invariant
of this situation is evidently the Postnikoff-Sinh invariant
\[ \alpha \in H^3(\pi_0, \pi_1) \ . \]

I have met these animals - without even looking for them -
in many situations, which I will not list now (I came across them
recently \uline{a} \uline{propos} the classification of ``ordinary'' formal groups
over a perfect field, in terms of \uline{affine} algebraic groups, or
\uline{commutative} formal groups, related by the strict Gr-structures of
this type (one relies on this formalism in an arbitrary topos
(not merely in (ENS)) - to make explicit the yoga that ``the
transcendent character of a formal group is concentrated essentially

\margin{131}

\clearpage
% R. Brown's translations of two letters to L. Breen, with his note and comments (cote 134-2): archive cote 134-2, site page 131 (PDF page 132).
% translation of the letter of 17-19 July 1975. Draft compared: ps1.tex, lines 1385-1424.
% AI-transcribed by Claude Opus 5.5 (claude-opus-5-5, adaptive thinking, effort high), Message Batches API, batch brown_full (msgbatch_01Md3i5yK1PxkpJn5GHfEExE), 2026-11-13; not yet checked by eye.
\scanpage{134-2/131}
\mspage{- 4 -}

in the commutative formal group'', discovered it seems by Dieudonné
...). The question which I wish to raise is the generalisation to
groups in (n-cat), where I expect to find a non-commutative chain
complex
\[ L_n \to L_{n-1} \to \dots \to L_1 \to L_0 \to 1 \]
with supplementary structures doubtless of the type of $\theta$ , but
what are they? It is agreed that the topological significance of
the structures is that they express exactly the ``truncated homotopy
type in $\dim \leq n$'' of topological groups, and even the homotopy
type in $\dim \leq n+1$ of connected topological spaces ...). Have
you candidates to propose?

[A commentary on biduality and homology is omitted.]

\mspage{\hand{132}}

\clearpage
% R. Brown's translations of two letters to L. Breen, with his note and comments (cote 134-2): archive cote 134-2, site page 132 (PDF page 133).
% translation of the letter of 17-19 July 1975. Draft compared: ps1.tex, lines 1538-1570.
% AI-transcribed by Claude Opus 5.5 (claude-opus-5-5, adaptive thinking, effort high), Message Batches API, batch brown_full (msgbatch_01Md3i5yK1PxkpJn5GHfEExE), 2026-11-13; not yet checked by eye.
\scanpage{134-2/132}
\mspage{- 5 -}

Your question 7 seems to indicate that there is discontent in
your mind on the significance of the ``homotopy type'' of $X$ , for
$X$ a topos (for examplethe étale topos of a scheme). Doubtless
you can understand the algebraic topology that one can do on $X$ ;
this includes semi-simplicial sheaves, the stacks of all kinds, and
the relations between these - from this point of view $X$ (with
its very rich structure of topos) virtually disappears so as to
become no more than a pale element of a ``homotopical category''
(or pro-homotopical), deduced from that of the topos by a process
of localisation
- at the first sight, all that still comes from this poor $X$ , are
the $\pi_i$ - and its cohomology groups with constant coefficients -
or at the worst twisted constants. When one digs more into this
definition of ``what is left to this poor $X$'' one falls precisely
on \uline{the} \uline{locally} \uline{constant} \uline{n-stacks}(as an $f : X \to X'$ which is a
homotopy equivalence induces a $(n+1)$-equivalence between the
mentioned locally constant n-stacks on $X$ and on $X'$). - which surely
contains the abelian chain complexes of length $n$ of sheaves with
locally constant cohomology, and the hyper-cohomology of these.
It is thus that one arrives at this triangle of objects which
mutually determine each other :

\begin{diagramblock}
\begin{tikzcd}[column sep=-3em, every arrow/.append style={<->}]
 & \begin{tabular}{@{}c@{}}
    topos (or topological space\\
    or semi-simplicial complex)\\
    \uline{modulo} n-homotopy
   \end{tabular}
   \arrow[dl]\arrow[dr] & \\
 \begin{tabular}{@{}l@{}}
    n-groupoids\\
    (up to n-equivalence)
   \end{tabular}
   \arrow[rr] &
 & \begin{tabular}{@{}l@{}}
    ``special'' n-categories\\
    (up to n-equivalence).
   \end{tabular}
\end{tikzcd}
\end{diagramblock}

\mspage{\hand{133}}

\clearpage
% R. Brown's translations of two letters to L. Breen, with his note and comments (cote 134-2): archive cote 134-2, site page 133 (PDF page 134).
% translation of the letter of 17-19 July 1975. Draft compared: ps1.tex, lines 1539-1570.
% AI-transcribed by Claude Opus 5.5 (claude-opus-5-5, adaptive thinking, effort high), Message Batches API, batch brown_full (msgbatch_01Md3i5yK1PxkpJn5GHfEExE), 2026-11-13; not yet checked by eye.
\scanpage{134-2/133}
\mspage{- 6 -}

One says that an n-category $E_n$ is ``special'' ( or \uline{n-galois}) if
it is n-equivalent to the category of locally constant (n - 1)-stacks
on a convenient topological space (or a topos), or what also turns
out, it is n-equivalent to the category of n-functors
$G_n \to (\text{n - Cat})$ , where $G_n$ is an n-groupoid. If $X$ , $G_n$ , $E_n$
correspond in this way, one calls $G_n$ the fundamental n-groupoid
of $X$ , or of $E_n$ , or that $E_n$ is the category of local n-systems
on $X$ , or on $G_n$ , or that $X$ is the geometric realisation of
$G_n$ or of $E_n$ , In analogy with the familiar case $n = 1$, $G_n$
can be interpreted as the full sub-n-category of $\underline{\mathrm{Hom}}_n(E_n,(\text{n - cat}))$
formed by the n-functors $E_n \to (\text{n-cat})$ satisfying certain
exactness properties (or what is equivalent: which commute with
finite $\varprojlim$ and arbitrary $\varinjlim$ ); but this raises the disquieting
vision of n-limits in n-categories. (N.B. The case $n = 2$ begins
to become familiar to us ...). It is prudent in all of this to
suppose that $X$ is ``locally homotopically trivial'', which ensures
the pro-simplicial set which Artin-Mazur associate to it (with
the help of nerves of hyper coverings) is essentially constant in
the ordinary homotopy category - then $X$ defines the usual
homotopy type. This is surely \uline{not} the case for the étale topos
of a scheme. In such case, the fundamental n-groupoid should be
conceived as a \uuline{pro}-\uline{n}-\uline{groupoid} (nothing surprising in that, in view
of the familiar theory of $\pi_1$), and $E_n$ as an (ind)-n-category
(the ind-structure corresponds to the exigencies of local triviality
for a variable n-stack, relative to coverings more and more fine on
X).

\margin{134}

\clearpage
% R. Brown's translations of two letters to L. Breen, with his note and comments (cote 134-2): archive cote 134-2, site page 134 (PDF page 135).
% translation of the letter of 17-19 July 1975. Draft compared: ps1.tex, lines 1563-1572.
% AI-transcribed by Claude Opus 5.5 (claude-opus-5-5, adaptive thinking, effort high), Message Batches API, batch brown_full (msgbatch_01Md3i5yK1PxkpJn5GHfEExE), 2026-11-13; not yet checked by eye.
\scanpage{134-2/134}
\mspage{- 7 -}

I nevertheless understand your instinctive resistance to
conceive this extreme stripping of a beautiful topos $X$ , to the
point of retaining only the meagre homotopy type. Even more I
am persuaded that going to the root of this instinctive resistance,
one arrives at a generalisation and deepening of the notion of
``homotopy type'', which bring us to still new waters in the
development of a good homotopical yoga. Here is what I wish to
come to pass.

Let us speak first of sheaves (of sets, or of modules, etc.)
instead of stacks, for simplicity, and place ourselves in the
étale topos of a scheme. The locally constant sheaves - modulo
a supplementary condition of finiteness which is sufficiently
anodyne - form the easiest of the \uline{constructible} sheaves, for
the definition of which they serve as models. Supposing $X$
quasi-compact and quasi-separated, the general constructible
sheaves are those for which there exist a finite partition
$X = \hand{\bigcup}_{i\in I} X_i$ of $X$ in ``cells'' or ``strata'' $X_i$ , each locally
closed and constructible, such that the restriction of $F$ to
every $X_i$ is locally constant (also a finiteness condition ...).
Thus the category of constructible sheaves on $X$ (which gives
back the category of all sheaves on passing to a category of
ind-objects ...) appears itself to be an inductive limit of
categories associated to finer and finer partitions on $X$ . One
can then, by such a finite partition $P$ , propose a more profound
study of the category of sheaves (or complexes of sheaves, or
stacks) which are ``P-constructible'' (or, more generally, which are

\mspage{\hand{135}}

\clearpage
% R. Brown's translations of two letters to L. Breen, with his note and comments (cote 134-2): archive cote 134-2, site page 135 (PDF page 136).
% translation of the letter of 17-19 July 1975. Draft compared: ps1.tex, lines 1570-1592.
% AI-transcribed by Claude Opus 5.5 (claude-opus-5-5, adaptive thinking, effort high), Message Batches API, batch brown_full (msgbatch_01Md3i5yK1PxkpJn5GHfEExE), 2026-11-13; not yet checked by eye.
\scanpage{134-2/135}
\mspage{- 8 -}

``locally constant'' on every $X_i$). These categories will not have
truly satisfying structures unless they are stable for standard
operations - such as $R\,\underline{\mathrm{Hom}}$ , or $Rj_*j^*$ where $j : X_i \to X$ is
a ``cell'' of the partition, etc. In fact, if $X$ is excellent and
one has resolution of singularities at ones disposal, one
knows that the torsion constructible sheaves (under the proviso
of being prime to the characteristic) are stable for all the
operations - but not only for a finite partition of $X$ but for
all. To have such a finer stability, it is necessary to make
some very strict hypotheses of ``equi-singularity'' on the given
stratification of $X$ , along the strata. I think nonetheless that
a sufficiency of known techniques show that $X$ admits arbitrarily
fine stratifications having the properties of equi-singularity
(and with the $X_i$ regular and connected, but this does not matter
for our present purposes). By way of example, suppose that there
are two closed strata $X_0$ and $X_1 = X\backslash X_0$ . After Artin's
devissage, the sheaf $F$ on $X$ comes from giving a sheaf
$F_1(= i_1^*\,F)$ on $X_1$ , and a homomorphism $F_0 \to i_0^*\,i_{1*}^*(F_1)$ , where
$i_0$ , $i_1$ are the inclusions $X_0 \xrightarrow{i_0} X \xrightarrow{i_1} X_1$ .
In order that $F$ should be $P$-constructible, it is necessary and
sufficient that $F_0$ and $F_1$ should be locally constant (plus
some accessory finiteness conditions...), on $X_0$ and $X_1$
respectively. Then (by virtue of the hypothesis of equi-singularity)
the same will be true of $X_1$ and the category of sheaves in which
we are interested can be expressed entirely in terms of the
category of locally constant sheaves on $X_0$ and $X_1$ , i.e. of
the mere homotopy type of $X_0$ and $X_1$ , except for the proviso

\margin{136}

\clearpage
% R. Brown's translations of two letters to L. Breen, with his note and comments (cote 134-2): archive cote 134-2, site page 136 (PDF page 137).
% translation of the letter of 17-19 July 1975. Draft compared: ps1.tex, lines 1570-1585.
% AI-transcribed by Claude Opus 5.5 (claude-opus-5-5, adaptive thinking, effort high), Message Batches API, batch brown_full (msgbatch_01Md3i5yK1PxkpJn5GHfEExE), 2026-11-13; not yet checked by eye.
\scanpage{134-2/136}
\mspage{- 9 -}

that we must make more explicit the nature of the left exact
functor $\phi$ . I think that this should be possible, in the context
of \uline{schemes} in which I am placed (technically rather sophisticated),
on introducing an ``étale tubular neighbourhood'' of $X_0$ in $X_1$
(which is a very satisfying topos, but not associated to a sheaf).
But this technical construction is only a paraphrase of a topol-%
ogical intuition extraordinarily simple, which one can make
explicit, on supposing to fix the ideas that the base field is
$\mathbb{C}$ and so one profits by working with locally compact spaces.
The topological idea behind the hypothesis of equi-singularity is
that there exist a \uline{tubular} \uline{neighbourhood} $T$ of $X_0$ in $X$
retracting onto $X_0$ and such that the pair $(X_0, T)$ over $X_0$
should be a locally trivial bundle, i.e. that $T\backslash X_0$ is locally
trivial over $X_0$ . In fact if $\partial T$ is the ``boundary'' of $T$ ,
which also should be a locally trivial bundle on $X_0$ , then $T/X$
is the bundle with fibres
$(\simeq (\partial T \times I) \amalg_{\partial T} X_0$ , where $I = [0, 1]$ , $\partial T \to \partial T \times I$ . is defined
by $x \longmapsto (x, 1)$ , and $\partial T \to X_0$ is the projection) then
$\mathring{T} = T\backslash X_0 \simeq \partial T \times [0, 1]$ is $X_0$-homotopic to $\partial T$ . If $X_0$ and
$X_1$ are non singular, then the same is true for $\mathring{T}$ and $\partial T$ ,
which are then topologically smooth fibrations on $X_0$ . Moreover,
putting $\tilde{X}_1 = X_1\backslash\mathring{T}$ , the inclusion $\tilde{X}_1 \to X_1$ is a homotopy equivalence
and $X$ can be recovered, up to homeomorphism, from the diagram
of spaces
\begin{diagramblock}\begin{tikzcd}
(\mathring{T} = T\backslash X_0 \cong)\partial T \ar[r, "j", "\text{inclusion}"'] \ar[d, "\text{fibration}"', "p"] & \tilde{X}_1(\cong X_1) \\
X_0 &
\end{tikzcd}\end{diagramblock}

\margin{137}

\clearpage
% R. Brown's translations of two letters to L. Breen, with his note and comments (cote 134-2): archive cote 134-2, site page 137 (PDF page 138).
% translation of the letter of 17-19 July 1975. Draft compared: ps1.tex, lines 1572-1612.
% AI-transcribed by Claude Opus 5.5 (claude-opus-5-5, adaptive thinking, effort high), Message Batches API, batch brown_full (msgbatch_01Md3i5yK1PxkpJn5GHfEExE), 2026-11-13; not yet checked by eye.
\scanpage{134-2/137}
\mspage{- 10 -}

as an amalgamated sum.  In terms of this diagram of spaces,
the above functor $\phi$ interprets immediately as
\[ \phi(F_1) \simeq p_* j^*(\tilde{F}_1) \]
where $F_1 \to \tilde{F}_1$ is the restriction from $X_1$ to $\tilde{X}_1$ (which is
an equivalence of categories for the envisaged (locally constant)
sheaves).  Giving $F = (F_0 , F_1 , u : F_0 \to \phi F_1)$ can then also
be made explicit as giving
\[ F_0 , \tilde{F}_1 , \tilde{u} : p^*(F_0) \to j^*(\tilde{F}_1) \]
where $F_0(\tilde{F}_1)$ are locally constant sheaves on $X_0$ (respectively
$X_1$).  It is necessary to recall that up to here $p$ is a real
fibration, and $j$ is an inclusion (in practice, - for the case
$X_0$ , $X_1$ smooth - the inclusion of the boundary is a manifold
with boundary).

If you prefer, one can also take the diagram which is less
pretty (but a little more canonical)
\begin{diagramblock}\begin{tikzcd}
\mathring{T} \ar[r,"j'"] \ar[d,swap,"p'"] & X_1 \\
X_0 &
\end{tikzcd}\end{diagramblock}
coming essentially to the same thing, as formed from spaces
homotopic to the preceding one.  One can even replace $X_0$ by $T$
($X_0$ being a deformation retract of it) and write
\begin{diagramblock}\begin{tikzcd}
\mathring{T} \ar[r] \ar[d,swap,"p''"] & X_1 \\
T &
\end{tikzcd}\end{diagramblock}

\margin{138}

\clearpage
% R. Brown's translations of two letters to L. Breen, with his note and comments (cote 134-2): archive cote 134-2, site page 138 (PDF page 139).
% translation of the letter of 17-19 July 1975. Draft compared: ps1.tex, lines 1585-1612.
% AI-transcribed by Claude Opus 5.5 (claude-opus-5-5, adaptive thinking, effort high), Message Batches API, batch brown_full (msgbatch_01Md3i5yK1PxkpJn5GHfEExE), 2026-11-13; not yet checked by eye.
\scanpage{134-2/138}
\mspage{- 11 -}

where ``literally'' $p''$ is now an inclusion, but ``morally'', it is
a \uline{fibration} with a very pretty fibre [notably in finite dimensions,
these are non-singular varieties - this is then much better than
that which is given by the yoga of Cartan-Serre ``every continuous
mapping is equivalent to a fibration''...] . This last diagram
is acceptable to a direct construction, purely algebraic, in the
context of schemes, once one has developed the construction of
étale tubular neighbourhood.

The point I wish to come to, is that the consideration of
sheaves (or complexes of stacks, of n-stacks ...) which are $P$-%
constructible on an $X$, where $P$ is a given ``equi singular''
stratification, comes in actual particular cases to the knowledge
of a diagram of ordinary \uline{homotopy} \uline{types} (or pro-types, if one
comes back to the étale topology)
\begin{diagramblock}\begin{tikzcd}
  \Delta \ar[r,"j"]\ar[d,swap,"p"] & X_1 \\
  X_0 &
\end{tikzcd}\end{diagramblock}
and taking local coefficients systems (or locally constant n-stacks)
on the vertices $X_0$ , $X_1$ , related by homeomorphisms of compat-%
ibility of the type $p^*(F_0) \to j^*(F_1)$ . It should be an amusing
excercise (which I have not yet done) to verify and to make
explicit how the usual operations on sheaves (on $X$ , or on a union
of strata of $X$ ) can be expressed in this dictionary, in the
stated case of non-singular strata (at this moment, there will be
a difficulty with the dualising complexes,

\mspage{\hand{139}}

\clearpage
% R. Brown's translations of two letters to L. Breen, with his note and comments (cote 134-2): archive cote 134-2, site page 139 (PDF page 140).
% translation of the letter of 17-19 July 1975. Draft compared: ps1.tex, lines 1605-1614.
% AI-transcribed by Claude Opus 5.5 (claude-opus-5-5, adaptive thinking, effort high), Message Batches API, batch brown_full (msgbatch_01Md3i5yK1PxkpJn5GHfEExE), 2026-11-13; not yet checked by eye.
\scanpage{134-2/139}
\mspage{- 12 -}

and with verifying the known formulas for these operations).
But it appears probable that, to carry out this transcription well,
it would be suitable, to regard a diagram of type
\begin{diagramblock}\begin{tikzcd}
  \cdot \ar[r]\ar[d] & {} \\
  {} &
\end{tikzcd}\end{diagramblock}
as based in the homotopy category formed from the category of
semi-simplicial sets, i.e. to consider the category of diagrams
of semi-simplicial sets, and to form from this a homotopy category
of functions.

I have recently more or less made explicit, while thinking
on \unsure{tne|the} foundations of ``tame topology'',(i.e. where one immediately
eliminates all wild phenomena) how an equisingular stratification,
say with non singular strata, of a compact ``tame  space'', gives
rise canonically to a diagram of spaces which are manifolds with
boundary, the arrows of the diagram being essentially locally
trivial fibrations of manifolds with boundary on the others (with
fibres manifolds with boundary), \uline{and} the inclusion of the boundary
in these manifolds with boundary [in fact, one finds slightly
more general inclusions, certain boundaries which appear being
provided with ``elementary'' cellular decomposition, i.e. the
closed strata are like manifolds with boundary which paste along
some of the common part of the boundary; and there is
% NOTE: the typist left the left half of the next line blank (a gap in the typed text, nothing hidden)
\hspace*{20em} the inclusions of these pieces
one in another ...], and can be reconstituted from this diagram
by pasting.  In other words, one has a canonical devissage
description of a tame compact space  $X$ , eventually union of an
equisingular stratification with non-singular strata, in terms
of finite diagrams of a precise nature made out of manifolds

\mspage{\hand{140}}

\clearpage
% R. Brown's translations of two letters to L. Breen, with his note and comments (cote 134-2): archive cote 134-2, site page 140 (PDF page 141).
% translation of the letter of 17-19 July 1975. Draft compared: ps1.tex, lines 1597-1614.
% AI-transcribed by Claude Opus 5.5 (claude-opus-5-5, adaptive thinking, effort high), Message Batches API, batch brown_full (msgbatch_01Md3i5yK1PxkpJn5GHfEExE), 2026-11-13; not yet checked by eye.
\scanpage{134-2/140}
\mspage{- 13 -}

with boundary.  When in a situation of sheaves (or complexes
of sheaves or n-stacks) which are $P$-constructible on $X$ , where
$P$ is such a fixed stratification, these also can be made
explicit in terms of the envisaged diagram, from which there is
no way to retrieve more than the ``homotopy type''. One thought
that the six operations on sheaves can be translated ad hoc to
this homotopical framework. Finally, if instead of having only
one compact tame space $X$ , one has a tame morphism $f : X \to Y$
then by choosing equisingular stratifications on $X$ and $Y$
adapted to $f$ (the strata of $X$ being notably locally trivial
fibrations on those of $Y$ ...), one should find a ``morphism''
of diagrams of manifolds with boundary expressing $X$ in that
expressing $Y$ (with some natural morphisms essentially equivalent
to fibrations of manifolds with boundaries on the others) in
such a way that the four operations $Rf_*$, $Rf_!$, $Lg^*$, $g^!$ between
$P_X$- and $P_Y$- constructible sheaves on $X$ and $Y$ (on locally
closed sub spaces $X'$ , $Y'$ which are union of strata, such that
$f$ induces $g : X' \to Y'$) can be expressed also in terms of
standard operations between homotopy type. Finally, all these
constructions, still partially hypothetical (there is work on
the foundations on the drawing board!) should be able to be
paraphrased in the framework of excellent schemes, using the
machinery of non-étale tubular neighbourhoods. In both cases,
the ``fine homotopy type'' of a tame space, respectively of an
excellent scheme, is defined by passage to the limit from ``$P$-
homotopy type'' associated to finer and finer equisingular
stratifications $P$ (with non-singular strata). This ``finer''
homotopy type would necessitate the knowledge, not only of sheaves
or locally constant n-stacks, but (via a passage to the inductive

\margin{141}

\clearpage
% R. Brown's translations of two letters to L. Breen, with his note and comments (cote 134-2): archive cote 134-2, site page 141 (PDF page 142).
% translation of the letter of 17-19 July 1975. Draft compared: ps1.tex, lines 1612-1621.
% AI-transcribed by Claude Opus 5.5 (claude-opus-5-5, adaptive thinking, effort high), Message Batches API, batch brown_full (msgbatch_01Md3i5yK1PxkpJn5GHfEExE), 2026-11-13; not yet checked by eye.
\scanpage{134-2/141}
\mspage{-14-}

limit) the knowledge of \uline{everything}.  And it depends in a convenient
sense, functorially on $X$ .  In the case of a scheme of finite
type on an algebraically closed  field $k$  say, the strongest
cohomological \uline{finiteness} \uline{theorem} in homotopy are expressed truly
precisely in terms of fine homotopy type, and say that the
ordinary homotopy types which are their elementary constituents
are essentially ``finite polyhedra'' - and even compact manifolds
with boundary - or in more precise fashion, the profinite
completions (in the sense of Artin-Mazur) prime to $p$ (the
characteristic) of such a space.  One sees clearly how to carry
out such a programme in characteristic  $0$ , but one has  more
fun, even    mystery , in the case  $p > 0$ , for the varieties
which, even  birationally, do not lift to characteristic $0$ !

From these essentially geometric thoughts, I have not at
this moment drawn up a precise programme for developing an
adequate algebraic structure to express them.  I restrict myself
to several marginal remarks.

For a long time I  have been intrigued by the idea of a
``linearisation'' of an (ordinary) homotopy type, i.e. questions of
the type: if  $X$  is a homotopy type, how much cohomological
information of the type: cohomology of  $X$  with variable co-%
efficients  $M$ (constant or twisted constant), multiplicative
structure  $H^{i}(X, M) \times H^{j}(X, N) \to H^{i+j}(X, M \otimes N)$ ,  then eventually
other cohomology operations - is it necessary to have to reconstruct
entirely the homotopy type (say, in this first setting, of the
type of derived pre-category, with the knowledge of the fundamental

\margin{142}

\clearpage
% R. Brown's translations of two letters to L. Breen, with his note and comments (cote 134-2): archive cote 134-2, site page 142 (PDF page 143).
% translation of the letter of 17-19 July 1975. Draft compared: ps1.tex, lines 1621-1623.
% AI-transcribed by Claude Opus 5.5 (claude-opus-5-5, adaptive thinking, effort high), Message Batches API, batch brown_full (msgbatch_01Md3i5yK1PxkpJn5GHfEExE), 2026-11-13; not yet checked by eye.
\scanpage{134-2/142}
\mspage{- 15 -}

group $\pi_1$ , the category of constant twisted coefficients
(= $\pi_1$-module), the functors $H^i(-, M)$ on these, with the structure
of cohomological functors relative to exact sequences, the
structure of \ \ product , etc. - related by certain formal
properties)? Once one has at ones disposal the language of derived
categories, the sub-category of the derived category of abelian
complexes, formed from complexes of sheaves of cohomology which
are locally constant on $X$ , with its triangulated structure
and its multiplicative structure $\overset{\mathbb{L}}{\otimes}$ (and eventually $\mathbb{R}\,\mathrm{Hom}$ ...)
give a more satisfying candidate for hoping to recover the
homotopy type. I have little idea if this suffices to permit
the recovery without effort, but on the contrary have no doubt
that on getting to the bottom of the ``linearisation'', that is
to say on working with \uline{non-abelian}, and with the (n+1)-categories
(without the supplementary structure on it!) of n-stacks of
constructible sheaves on $X$ , for all $n$ , one arrives at a
reconstruction of the homotopy type via its fundamental $\infty$-groupoid,
as explained in my previous letter and recalled in this. (This
signifies in particular that all the possible and imaginable
cohomology operations are already induced in giving such a system
of n-categories ...).

In the same way, the more complicated homotopy type, related
to certain finite diagrams, which one can associate to certain
types of stratification $P$ of tame topological spaces $X$ , let's
say, should correspond in an also perfect fashion to the (n+1)-%
category of n-stacks on $X$ which are locally constant on each
of the strata of $P$ (say: which are subordinated to $P$).
If the above programme for the locally constant derived category

\mspage{\hand{143}}

\clearpage
% R. Brown's translations of two letters to L. Breen, with his note and comments (cote 134-2): archive cote 134-2, site page 143 (PDF page 144).
% translation of the letter of 17-19 July 1975. Draft compared: ps1.tex, lines 1621-1647.
% AI-transcribed by Claude Opus 5.5 (claude-opus-5-5, adaptive thinking, effort high), Message Batches API, batch brown_full (msgbatch_01Md3i5yK1PxkpJn5GHfEExE), 2026-11-13; not yet checked by eye.
\scanpage{134-2/143}
\mspage{- 16 -}

works beautifully and well , one expects here to recover the
mixed homotopy type from the sub-category  of the derived category
of abelian sheaves on  $X$ , made from the complexes which are
cohomologically locally constant on all of  the strata - with
also the operations $\overset{\mathbb{L}}{\otimes}$ , $\mathbb{R}\mathrm{Hom}$, and, evidently, the four operations
$\mathbb{R}g_!$ , $\mathbb{R}g_*$ , $Lg^*$ , $g^!$  for the inclusions  $g : Z' \to Z''$ of the
various locally closed unions of strata ... .  The problem here
is that we don't yet even know what is a triangulated category,
nor even what is its non-commutative version, described probably
more simply and more fundamentally, as a ``homotopical category''
with operations of ``fibrations'' and ``cofibrations''.  [N.B. A
first approach has been developed by Quillen at the beginning
of the years ``70''.]

It is surely time that I finish this ``lettre-fleuve'', which
is becoming more and more verbose.  Just one question: what is
this marvellous formula of Bloch - Quillen to which you made
allusion, and of which I have never heard speak, and which makes
my mouth water?

Very cordially to you,

Alexandre  Grothendieck.

\mspage{\hand{144}}

\end{document}
